\documentclass[11pt]{article}
\usepackage[margin=0.75in]{geometry}
\usepackage{times}
\usepackage{booktabs}
\usepackage{graphicx}
\usepackage{hyperref}
\usepackage{parskip}
\usepackage{enumitem}
\usepackage{caption}
\usepackage{float}

\setlength{\parskip}{2pt}
\setlength{\parindent}{0pt}

\hypersetup{colorlinks=true, linkcolor=black, citecolor=black, urlcolor=blue}

\title{\vspace{-1.5em}\textbf{Don't Trust All Your Neighbors: Re-implementing Graph Attention Networks}\vspace{-0.5em}}
\author{Yash Chatha (yc2727) \quad Jerry Ji (rj378) \quad Leon Huang (lyh7) \\
\small CS 4782: Deep Learning, Cornell University, Spring 2026}
\date{\vspace{-1em}}

\begin{document}
\maketitle
\vspace{-1em}

\section*{1. Introduction}

Graph neural networks address node classification on graph-structured data by aggregating information from a node's local neighborhood. This project re-implements \textbf{Graph Attention Networks (GAT)} by Veli\v{c}kovi\'{c} et al. (ICLR 2018). The paper's primary contribution is replacing the fixed equal-weight neighbor aggregation of Graph Convolutional Networks (GCN) with a learned attention mechanism: for each node, the model computes a separate attention score per edge, weighting neighbors by relevance rather than treating them uniformly. This is achieved without requiring knowledge of the full graph structure upfront.

\section*{2. Chosen Result}

We reproduce \textbf{Table 2}: transductive node classification accuracy on Cora and CiteSeer averaged over 100 runs. The paper reports $83.0 \pm 0.7\%$ on Cora and $72.5 \pm 0.7\%$ on CiteSeer; this is the central claim that learned attention outperforms fixed-weight GCN aggregation ($81.5\%$ / $70.3\%$). Reproducing it validates the architecture and training protocol, and establishes a baseline for our extensions.

\section*{3. Methodology}

\textbf{Architecture.} We implement a 2-layer GAT using PyTorch Geometric's \texttt{GATConv}. Layer 1 uses $K=8$ attention heads each computing $F'=8$ features, concatenated to a 64-dimensional node representation with ELU activation. Layer 2 uses a single attention head outputting $C$ class logits followed by log-softmax. Dropout $p=0.6$ is applied to both layer inputs and attention coefficients, matching the paper's regularization exactly.

\textbf{Training.} Adam optimizer, $\text{lr}=0.005$, L2 weight decay $\lambda=0.0005$, up to 10,000 epochs with early stopping (patience 100). \textbf{Modification:} we monitor only validation loss for early stopping, rather than both validation loss and accuracy as in the original TensorFlow implementation. This was a deliberate choice to ensure a consistent, comparable training protocol across all experiments, particularly the low-label benchmark where label scarcity introduces noise that makes accuracy-based stopping less stable.

\textbf{Datasets.} Cora (2,708 nodes, 1,433 features, 7 classes) and CiteSeer (3,327 nodes, 3,703 features, 6 classes), loaded via PyG Planetoid with NormalizeFeatures. Standard public splits: 20 nodes/class train, 500 val, 1,000 test.

\textbf{Evaluation.} 100 seeds; mean $\pm$ std test accuracy reported.

\textbf{Extensions.} Beyond reproduction, we designed three independent experiments: (1) a \textit{low-label benchmark} varying the fraction of training labels from 10\% to 100\%; (2) \textit{TinyGAT distillation}, training a 6$\times$ smaller student model using the full GAT as a knowledge distillation teacher; and (3) a \textit{failure case explorer}, analyzing every misclassified test node's confidence, neighbor class distribution, and per-edge attention weights.

\section*{4. Results \& Analysis}

\textbf{Baseline reproduction.}

\begin{table}[H]
\centering
\small
\begin{tabular}{lccc}
\toprule
Dataset & Paper & Ours & Gap \\
\midrule
Cora & $83.0 \pm 0.7\%$ & $\mathbf{83.22 \pm 0.41\%}$ & $+0.22\%$ \\
CiteSeer & $72.5 \pm 0.7\%$ & $70.94 \pm 0.52\%$ & $-1.56\%$ \\
\bottomrule
\end{tabular}
\end{table}

Cora matches and slightly exceeds the paper's target. The CiteSeer gap is attributable to our early stopping deviation: monitoring loss only is a less sensitive stopping signal than loss-plus-accuracy, which may cause the model to stop at a suboptimal point on CiteSeer's harder optimization landscape.

\textbf{Low-label benchmark.} Accuracy degrades monotonically as labels decrease, but the magnitude differs sharply. On Cora, performance at 10\% labels (14 nodes) is $58.8 \pm 5.9\%$, well above random chance ($14.3\%$). On CiteSeer, the same fraction (12 nodes) collapses to $30.2 \pm 17.9\%$, with one seed hitting $8.3\%$, below random chance ($16.7\%$). We attribute this to CiteSeer's sparser feature vocabulary (3,703-dim with few non-zero entries per node), leaving the model with negligible discriminative signal at extreme label scarcity.

\textbf{TinyGAT distillation.} TinyGAT (2 heads $\times$ 4 features, $\sim$6$\times$ fewer parameters) trained with KL-divergence soft targets (temperature=2.0, $\alpha=0.7$) showed mean distillation gain of $-0.24\%$ on Cora and $-0.66\%$ on CiteSeer. TinyGAT with supervision alone already achieves $82.5\%$ / $70.8\%$, nearly matching the teacher. Model capacity is not the bottleneck; both models saturate the available signal.

\textbf{Failure case analysis.} We analyzed 25 Cora and 24 CiteSeer misclassifications for seed 1. Two failure modes emerged: (1) \textit{ambiguous neighborhoods}, where nodes with mixed-class neighbors receive nearly uniform attention weights, producing low-confidence errors that may be irreducible; and (2) \textit{hub node influence}, where a single high-attention wrong-class neighbor dominates aggregation. On Cora, node 1358 (class 2) appears in the top-5 attention sources for 6 of 25 failures, with two errors exceeding 92\% confidence. Since attention weights are lower-bounded at zero (post-softmax), the model can reduce but never cancel a misleading neighbor's contribution.

\section*{5. Reflections}

\textbf{Key lessons.} Hyperparameter choices interact non-trivially with dataset difficulty: our early stopping deviation (loss-only) costs approximately 1.5\% on CiteSeer but is negligible on Cora. Future re-implementations should match the exact stopping criterion before declaring a gap.

\textbf{Surprising findings.} CiteSeer's collapse to near-random at 10\% labels was unexpected; standard benchmark results do not expose this feature-sparsity sensitivity. The distillation null result reframes Cora and CiteSeer as capacity-saturated benchmarks where architecture above a low threshold doesn't matter.

\textbf{Future directions.} (1) Run low-label and distillation experiments on PubMed and PPI to test generalization at scale; (2) explore repulsive attention mechanisms that can down-weight misleading hub neighbors beyond what softmax allows; (3) ablate CiteSeer's low-label collapse by independently controlling feature density and graph topology.

\section*{References}

\begin{enumerate}[leftmargin=*, label={[\arabic*]}]
\item Veli\v{c}kovi\'{c}, P., Cucurull, G., Casanova, A., Romero, A., Li\`{o}, P., \& Bengio, Y. (2018). Graph Attention Networks. \textit{ICLR 2018}.
\item Fey, M. \& Lenssen, J.E. (2019). Fast Graph Representation Learning with PyTorch Geometric. \textit{ICLR Workshop on Representation Learning on Graphs and Manifolds}.
\item Yang, Z., Cohen, W., \& Salakhutdinov, R. (2016). Revisiting Semi-Supervised Learning with Graph Embeddings. \textit{ICML 2016}.
\item Kipf, T.N. \& Welling, M. (2017). Semi-Supervised Classification with Graph Convolutional Networks. \textit{ICLR 2017}.
\end{enumerate}

\end{document}
